\documentclass[10pt,a4paper]{amsart}
\usepackage[margin=19mm]{geometry}
\usepackage{amsmath,amssymb,mathtools}
\usepackage[scr=rsfs,cal=euler]{mathalfa}
\usepackage{xfrac}
\usepackage[unicode,hidelinks,bookmarks=false]{hyperref}
\newcommand{\ie}{i.e.\ }
\newcommand{\cf}{cf.\ }
\newcommand{\resp}{respectively\ }
\newcommand{\Subsection}{\subsection}
\providecommand{\nobreakdash}{\nobreak\hbox{-}}
\setlength{\emergencystretch}{2em}
\multlinegap0pt
\usepackage{lmodern}
\usepackage{mathtools}
\usepackage{enumitem}
\hypersetup{pdftitle={Sharp spectral norm concentration of sparse random tensors},
pdfauthor={Zhixin Zhou and Yizhe Zhu}}
\numberwithin{equation}{section}
\newtheorem{theorem}{Theorem}[section]
\newtheorem{lemma}[theorem]{Lemma}
\newtheorem{proposition}[theorem]{Proposition}
\newtheorem{corollary}[theorem]{Corollary}
\newtheorem{definition}[theorem]{Definition}
\newtheorem{remark}[theorem]{Remark}
\newcommand{\E}{\mathbb E}
\newcommand{\Pp}{\mathbb P}
\newcommand{\R}{\mathbb R}
\newcommand{\one}{\mathbf 1}
\newcommand{\eps}{\varepsilon}
\begin{document}
\title[Sharp spectral norm concentration of sparse random tensors]{Sharp spectral norm concentration of sparse random tensors}
\author{Zhixin Zhou, Yizhe Zhu}
\date{}
\maketitle
\noindent\emph{Source and licence.} Sharp spectral norm concentration of sparse random tensors. Version arXiv:2609.20520v1 (CC BY 4.0). This material is distributed under the Creative Commons Attribution 4.0 International licence, http://creativecommons.org/licenses/by/4.0/, as recorded in the arXiv OAI licence metadata for this exact version and shown on the abstract page https://arxiv.org/abs/2609.20520v1. Full rights notice: licenses/arxiv-2609.20520-CC-BY-4.0.txt.\par\smallskip
\noindent\textit{Editorial scope.} Counted unit: the original \texttt{theorem} environment \texttt{thm:main} at source lines 352--359 together with its complete proof at lines 1081--1148; the delivered document keeps source lines 214--1400 so that every referenced label and every citation resolves inside it. Native editable source is the paper's own concatenated TeX; the AMS wrapper is editorial formatting only. No publisher class, upstream script or unrelated artwork is redistributed.
\par\smallskip\noindent\textit{Reference note.} This article ships no compiled bibliography (biblatex); the entries delivered here are composed from the author's own BibTeX (.bib) fields, with no field invented and no entry dropped.
\begin{equation}\label{eq:homogeneity}
 |T(x_1,\ldots,x_k)|\le\|T\|\prod_{j=1}^k\|x_j\|_2.
\end{equation}


Denote
$\log_+u=\max\{0,\log u\}$. The letter $C$ denotes a positive
constant that may change from line to line; subscripts indicate its
allowed dependence. The order $k$ is fixed throughout.

\subsection{Related work}\label{sec:related}

The light--heavy method originates in the Kahn--Szemer\'edi approach
to random regular graphs~\cite{FKS89}. Feige and Ofek~\cite{FO05} applied spectral
and discrepancy methods to sparse random graphs, including the
logarithmic-degree regime and the effect of removing high-degree
vertices.
Subsequent developments include spectral estimates for a wide range of random graph
and hypergraph models~\cite{BFSU98,LR15,DJPP13,CGJ18,TY19,Z23,DWZ25}.
Le, Levina, and Vershynin~\cite{LLV17} obtained optimal-order matrix
concentration after degree regularization by combining
Grothendieck--Pietsch factorization with a decomposition of the random
graph.

For random tensors, Nguyen, Drineas, and Tran~\cite{NDT15} studied spectral
norms through entropy and concentration estimates for sparse and spread
vectors, with applications to tensor sparsification.
Tomioka and Suzuki~\cite{TS14} gave a covering-number argument under
sub-Gaussian assumptions on fixed multilinear evaluations.
For Bernoulli$(p)$ entries, the sub-Gaussian bound is not uniform at
the variance scale $p(1-p)$ as $p$ tends to zero.
For real and complex tensors with independent standard Gaussian
entries, Dartois and McKenna~\cite{DM24} obtained
high-probability upper bounds with refined leading constants by
counting critical points of the associated Gaussian multilinear form.
A complementary moment method bounds expected injective norms in
Gaussian and certain non-Gaussian models~\cite{DM26}.

Concentration for sampled tensors also arises in tensor completion
and network estimation. Jain and Oh~\cite{JO14}
proved a spectral-norm bound for randomly sampled symmetric
third-order tensors in their analysis of tensor completion.
Lei, Chen, and Lynch~\cite{LCL20} obtained a spectral-norm
concentration bound for the centered third-order adjacency tensor
in a multilayer network model.
For tensor sparsification, Xia and Yuan~\cite{XY21}
established a spectral-norm concentration inequality for tensors
obtained by entrywise Bernoulli sampling and reweighting.


Guan, Jiang, and Li~\cite{GJL25} observed that a log-free weighted Bernoulli
spectral-norm estimate would remove the remaining order-dependent
logarithmic factor from their tensor robust-PCA recovery conditions.
Our Corollary~\ref{cor:weighted}  therefore
proves Conjecture~7.1 in \cite{GJL25}.

Concentration also depends on the choice of tensor norm.
Yuan and Zhang~\cite{YZ17} proved concentration
inequalities for sampled tensors under an \emph{incoherent spectral
norm}. This norm restricts the sizes of the coordinates of the test
vectors in all but two modes, which allows them to show concentration below $np\asymp 1$.
A sharper sparse tensor concentration bound under this incoherent
norm, with an application to tensor completion, was obtained in~\cite{LMSZ26}.
These results restrict the test vectors, whereas the spectral norm
in this paper is taken over all Euclidean unit vectors.

For independent entries, Boedihardjo~\cite{Boe24} obtained
expected injective-norm bounds in terms of the entry variances,
first for Gaussian tensors and then for bounded entries, together
with concentration about the expectation.
In our homogeneous Bernoulli model, his Corollary~1.5 gives
\[
 \E\|T-\E T\|
 \le C_k\bigl(\sqrt{np(1-p)}+(\log n)^2\bigr).
\]
The additive logarithmic term does not yield the $\sqrt{np}$ scale
in the regime $np\asymp \log n$.

Lucca and Pesenti~\cite[Theorem~1.3]{LP26} determined the
spectral gap of sparse Erd\H{o}s--R\'enyi hypergraphs.
For fixed $k\ge3$, let $T$ be the symmetric $0$--$1$ adjacency
tensor of a $k$-uniform hypergraph on $n$ vertices, with each
$k$-element subset included independently with probability $p$.
If $pn^{k/2}\to\infty$ as $n\to\infty$, then with high probability,
\[
 \|T\|-\|T-\E T\|=(1-o(1))pn^{k/2}.
\]
This gives a spectral-gap asymptotic in a regime extending below
$p=1/n$, whereas our main result gives a sharp $O(\sqrt{np})$
centered-norm bound for independent-entry tensors when $np\ge c\log n$.

Another line of work concerns tensors with dependent entries.
For tensor products of independent random vectors, available results include
concentration for convex Lipschitz functions under bounded-coordinate
assumptions~\cite{V20}, as well as Hanson--Wright-type inequalities
and concentration of Euclidean norms of linear images when the underlying
coordinates are independent and sub-Gaussian~\cite{BKW22,V20}.
Nonasymptotic injective-norm bounds have also been established for Gaussian
and Rademacher tensor series with correlated entries~\cite{BGJLR24}.
For centered empirical moment tensors and their asymmetric counterparts,
sharp injective-norm bounds under sub-Gaussian assumptions are obtained
through empirical-process methods~\cite{ACS25,CS26}.
In these empirical models, the observations are independent, but the
tensor entries generally share random factors and are therefore dependent.
These results complement the independent-entry concentration problem
considered here.

Friedman and Wigderson~\cite{FW95}
studied the second eigenvalue of a random hypergraph generated by a
fixed number of independent uniform ordered edge draws, with replacement.
This differs from the independent-entry Bernoulli model. Our proof gives an improved bound on the second eigenvalue in their random hypergraph model.


\subsection*{Organization of the paper}
The paper is organized as follows. Section~\ref{sec:main} states
the Bernoulli concentration theorem and its weighted extension.
Section~\ref{sec:fw} gives the application to the
Friedman--Wigderson model. The proofs are given in
Sections~\ref{sec:proof} and~\ref{sec:fw-proof}, respectively.
Appendix~\ref{app:dyadic} contains the dyadic entropy estimate used
in the heavy-tuple argument.

\section{Main results}\label{sec:main}

Let $T$ have independent entries
\begin{equation}\label{eq:model}
 t_{\mathbf i}\sim\operatorname{Bernoulli}(p),
 \qquad \mathbf i\in[n]^k.
\end{equation}
Thus the mean tensor and the centered tensor are
\[
 P=\E T=p\,\one^{\otimes k},
 \qquad W=(w_{\mathbf i})=T-P.
\]
Independence in \eqref{eq:model} is over all ordered index tuples.
Set $d=np$; this is the expected number of nonzero entries in any
one-dimensional fiber.


\begin{theorem}\label{thm:main}
Let $k\ge2$ be a fixed integer. For any $c>0$ and $r>0$, there exists a
constant $C_{k,r,c}$ depending only on $k,r,c$ such that, for all
$n\ge2$ and $p\in(0,1]$, if $ d=np\ge c\log n$,  then the tensor in \eqref{eq:model} satisfies
\begin{equation}\label{eq:main}
 \Pp\{\|T-\E T\|>C_{k,r,c}\sqrt d\}\le n^{-r}.
\end{equation}
\end{theorem}


Here $c>0$ is arbitrary and fixed, and $C_{k,r,c}$ may depend on $c$,
$k$, and $r$.

\medskip
\noindent\textit{Optimality of the scale and sparsity condition.}
The tensor norm always dominates the largest $\ell_2$-norm of its
one-dimensional fibers. In particular,
\[
 \|T-\E T\|\ge\max_f\|f\|_2,
\]
where the maximum ranges over all one-dimensional fibers $f$ of
$T-\E T$. This follows by fixing all but one argument at standard
basis vectors. For $p\le1/2$ and $d\ge c\log n$, even a fixed fiber
has norm at least $\sqrt{d/2}$ with probability at least
$1-e^{-d/8}$ by the Chernoff bound, showing that the $\sqrt d$ scale
is optimal. If $1\le d=np=o(\log n)$, the maximum fiber norm is
$\omega(\sqrt d)$ with high probability, so
$\|T-\E T\|/\sqrt d\to\infty$ in probability as $n\to\infty$.
Thus the condition $d\ge c\log n$ is optimal in order.


\bigskip 

The upper bound in Theorem~\ref{thm:main} also holds for entrywise
Bernoulli sampling of an arbitrary
deterministic tensor. 

\begin{corollary}[Weighted Bernoulli sampling]\label{cor:weighted}
Let $k\ge2$ be fixed, let $1\le n_1,\ldots,n_k\le N$ with $N\ge2$,
and let $A=(a_{\mathbf i})\in\R^{n_1\times\cdots\times n_k}$ be
deterministic. Let $\Omega=(\omega_{\mathbf i})$ have independent
entries with
\[
 \omega_{i_1,\ldots,i_k}\sim
 \operatorname{Bernoulli}(p_{i_1,\ldots,i_k}),
 \qquad 0\le p_{i_1,\ldots,i_k}\le1.
\]
Write $p_{\mathbf i}=p_{i_1,\ldots,i_k}$ and put
$p_*=\max_{\mathbf i}p_{\mathbf i}$.
For any $c,r>0$, there is a constant $C_{k,r,c}$ such that, if
$Np_*\ge c\log N$, then
\begin{equation}\label{eq:weighted}
 \Pp\left\{\|A\odot(\Omega-\E\Omega)\|>
 C_{k,r,c}\|A\|_\infty\sqrt{Np_*}\right\}\le N^{-r},
\end{equation}
where $\odot$ denotes entrywise multiplication and
$\|A\|_\infty=\max_{\mathbf i}|a_{\mathbf i}|$.
\end{corollary}

\section{Application to the Friedman--Wigderson model}\label{sec:fw}

We also consider the random hypergraph model of Friedman and
Wigderson~\cite[Section~6]{FW95}. Here the number of sampled
edges is fixed. The edge draws are independent, but the resulting
tensor entries are dependent. We state the concentration bound
in this section and give its proof in Section~\ref{sec:fw-proof}.

Let $m\ge1$ be an integer, and let $E_1,\ldots,E_m$ be independent
uniform random elements of $[n]^k$. Repetitions of an edge and
repeated coordinates within an edge are allowed. Define the
occupancy tensor $X=(X_{\mathbf i})$ by
\begin{equation}\label{eq:fw-model}
 X_{\mathbf i}=\sum_{\nu=1}^m\one\{E_\nu=\mathbf i\},
 \qquad d=\frac{m}{n^{k-1}},\qquad
 P=\E X=\frac dn\one^{\otimes k}.
\end{equation}
In this section $P$ denotes the mean of $X$, so $\E X_{\mathbf i}=d/n$.
The Friedman--Wigderson second eigenvalue is
\begin{equation}\label{eq:fw-lambda}
 \lambda_{2,d}(X)
 =\|X-P\|.
\end{equation}
This definition specifies the centering, not an exact regularity
constraint on the sampled tensor. Every fiber has expected degree $d$,
but its realized degree is random.



\begin{theorem}[Log-free second eigenvalue bound]\label{thm:fw}
Let $k\ge2$ be a fixed integer. For any $c,r>0$, there exists a constant
$C_{k,r,c}$ depending only on $k,r,c$ such that, for all $n\ge2$ and
$m\ge1$, if $ d\ge c\log n$, the tensor in \eqref{eq:fw-model} satisfies
\[
 \Pp\{\lambda_{2,d}(X)>C_{k,r,c}\sqrt d\}\le n^{-r}.
\]
\end{theorem}

For fixed $k\ge3$ and the same  model,
Friedman and Wigderson~\cite[Theorem~1.5]{FW95} proved, for
$d\ge Ck\log n$, that with high probability
 $\lambda_{2,d}(X)\le (C\log n)^{k/2}\sqrt d$.
Thus Theorem~\ref{thm:fw} removes the factor $(\log n)^{k/2}$.




\section{Proofs for the Bernoulli model}\label{sec:proof}

We first discretize the unit sphere and estimate the light tuples.
For the heavy tuples, we establish the fiber-degree and discrepancy
bounds and then carry out a deterministic summation over dyadic
classes. Throughout this section, $T$ is the Bernoulli tensor in
\eqref{eq:model}, $P=\E T$, and $d=np$, except in the explicitly
deterministic statement and proof of Proposition~\ref{prop:heavy}.

\subsection{Discretization}\label{sec:discretization}

We use a Euclidean net on the unit sphere in $\R^n$.
A set is an $\eps$-net if every point of the sphere is within
Euclidean distance $\eps$ of one of its points.
Unlike the lattice discretization in \cite[Section~4.1]{ZZ21},
the net below is chosen directly on the sphere.
\begin{lemma}[Euclidean sphere net {\cite[Corollary~4.2.11]{V26}}]\label{lem:net}
For $0<\eps<1$, the Euclidean unit sphere in $\R^n$ has a deterministic
$\eps$-net with at most $(1+2/\eps)^n$ elements.
\end{lemma}

Fix $\eps=1/(2k)$ and choose such a deterministic net $\mathcal T$.
A product net consists of all $k$-tuples in $\mathcal T^k$;
it has at most $(1+4k)^{kn}$ elements.
For any tensor $W$, this net gives
\begin{equation}\label{eq:net-bound}
 \|W\|\le2\max_{(y_1,\ldots,y_k)\in\mathcal T^k}
                       |W(y_1,\ldots,y_k)|.
\end{equation}
Indeed, for unit vectors $x_j$, choose $y_j\in\mathcal T$ with
$\|x_j-y_j\|_2\le1/(2k)$. Replacing the vectors one at a time and
using \eqref{eq:homogeneity} gives
\[
 |W(x_1,\ldots,x_k)-W(y_1,\ldots,y_k)|
 \le\sum_{j=1}^k\|W\|\|x_j-y_j\|_2
 \le\frac12\|W\|.
\]
Taking the supremum proves \eqref{eq:net-bound}. We will apply this
inequality to the full multilinear form after combining the light
and heavy contributions.

Put $\tau=\sqrt d/n$. For $x=(x_1,\ldots,x_k)$, define the set of light tuples by
\begin{equation}\label{eq:light-set}
 \mathcal L(x)
 =\left\{\mathbf i\in[n]^k:
             \left|\prod_{j=1}^k x_{j,i_j}\right|\le\tau\right\}.
\end{equation}
The remaining tuples are heavy. The threshold is imposed on the
coordinate product. We will therefore
decompose all the vectors into coordinate classes when estimating
the heavy tuples.

\subsection{Light tuples}\label{sec:light}

\begin{lemma}[The light part on a product net]\label{lem:light}
Fix $k$ and $R>0$. Put $W=T-P$ and $\tau=\sqrt d/n$.
Let $\mathcal T$ be a deterministic $1/(2k)$-net of the Euclidean unit sphere,
with $|\mathcal T|\le(1+4k)^n$. For a tuple
$x=(x_1,\ldots,x_k)\in\mathcal T^k$, put
\[
 L_W(x)=\sum_{\mathbf i\in\mathcal L(x)}
           w_{\mathbf i}\prod_jx_{j,i_j}.
\]
There is a constant $B_{k,R}$ such that
\[
 \Pp\left\{\max_{x\in\mathcal T^k}|L_W(x)|>
                         B_{k,R}\sqrt d\right\}\le n^{-R}.
\]
\end{lemma}
\begin{proof}
For a fixed tuple, the summands are independent, centered, and bounded
in absolute value by $\tau$. Their total variance is at most
\[
 p\sum_{\mathbf i}\prod_jx_{j,i_j}^2
       =p\prod_j\|x_j\|_2^2=p.
\]
Bernstein's inequality~\cite[Theorem~2.9.5]{V26} gives
\[
 \Pp\{|L_W(x)|>B\sqrt d\}
 \le2\exp\left\{-\frac{B^2d}{2p+(2/3)\tau B\sqrt d}\right\}
 =2\exp\left\{-n\frac{B^2}{2+(2/3)B}\right\}.
\]
Choose $B=B_{k,R}$ so that
\[
 \frac{B^2}{2+(2/3)B}\ge k\log(1+4k)+R+1.
\]
A union bound over at most $(1+4k)^{kn}$ tuples then gives a failure
probability at most $2e^{-(R+1)n}\le n^{-R}$ for $n\ge2$.
\end{proof}


\subsection{Heavy tuples}\label{sec:heavy}

A \emph{box} is a set of the form
\[
 Q=S_1\times\cdots\times S_k,
 \qquad S_j\subseteq[n]\quad(j=1,\ldots,k).
\]
The set $S_j$ is the \emph{coordinate set} of $Q$ in mode $j$.

We estimate the heavy contribution of the nonnegative tensor $T$;
the corresponding contribution of its mean is at most $\sqrt d$,
as shown in \eqref{eq:heavy-mean}. We need two probabilistic inputs:
a bound on all fiber degrees and a joint discrepancy estimate for
collections of boxes. The rest of the argument is deterministic.
The auxiliary entropy estimate in Appendix~\ref{app:dyadic} is used
to sum over dyadic levels without a factor depending on their number.

\subsubsection{The bounded degree property}\label{sec:degrees}

A one-dimensional fiber is obtained by fixing all indices except one.
For a fiber $F\subseteq[n]^k$, its degree is simply
\[
 \deg(F)=\sum_{\mathbf i\in F}t_{\mathbf i}.
\]
There are at most $kn^{k-1}$ such fibers, and each degree has mean $d$.
These are degrees of $(k-1)$-tuples, not vertex degrees of a
$k$-uniform hypergraph.
\begin{lemma}[Maximum fiber degree]\label{lem:fibers}
For fixed $k\ge2$, $R>0$, and $c>0$, there is a constant
$\kappa=\kappa(k,R,c)>1$ such that, if $d=np\ge c\log n$, all
one-dimensional fibers of $T$ have degree at most $\kappa d$,
except on an event of probability at most
$n^{-R}$.
\end{lemma}
\begin{proof}
There are at most $kn^{k-1}$ fibers and each degree is
$\operatorname{Bin}(n,p)$. Chernoff's bound gives, for $\kappa>1$,
\[
 \Pp\{\operatorname{Bin}(n,p)\ge\kappa d\}
 \le\exp\{-d(\kappa\log\kappa-\kappa+1)\}.
\]
Choose $\kappa$ such that
$\kappa\log\kappa-\kappa+1\ge (R+k+1+\log_2 k)/c$. The union bound and
$d\ge c\log n$ then give a failure probability at most
\[
 kn^{k-1}n^{-(R+k+1+\log_2 k)}
 =n^{-R-2}\,k n^{-\log_2 k}\le n^{-R}.
\]
\end{proof}

\subsubsection{Simultaneous multibox discrepancy}\label{sec:entropy}

We measure the discrepancy of a box by an upper-tail rate. For a
mean $\mu$ and a count $u\gg\mu$, the Chernoff exponent has order
$u\log(u/\mu)$. Counts in disjoint boxes are independent, so an
exponential-moment bound for one rate gives a bound for the sum of
the rates over a collection of boxes.

For $\mu>0$ and $u\ge0$, define the upper-tail Poisson rate
\begin{equation}\label{eq:rate}
 I_\mu(u)=
 \begin{cases}
 u\log(u/\mu)-u+\mu,&u\ge\mu,\\
 0,&0\le u<\mu.
 \end{cases}
\end{equation}
For a binomial count $M$ with mean $\mu$, the classical Chernoff
bound~\cite[Theorem~2.3.1]{V26} gives
\[
 \Pp\{M\ge u\}\le e^{-I_\mu(u)},\qquad u\ge\mu.
\]
Thus a large value of $I_\mu(u)$ means
that observing a count at least $u$ has small probability.
For $u<\mu$, we set $I_\mu(u)=0$ so that counts below the mean
contribute nothing to the upper-tail discrepancy. In this range,
$e^{-I_\mu(u)}=1$ gives only the trivial probability bound
$\Pp\{M\ge u\}\le1$.

\begin{lemma}[Exponential moment of the rate]\label{lem:rate}
If $M\sim\operatorname{Bin}(V,p)$, where $V\ge1$ and $\mu=pV>0$, then
\[
 \Pp\{I_\mu(M)>t\}\le e^{-t}\quad(t\ge0),
 \qquad
 \E\exp\bigl(I_\mu(M)/2\bigr)\le2.
\]
\end{lemma}
\begin{proof}
For each $t\ge0$, let $u_t\ge\mu$ be the unique solution of
$I_\mu(u_t)=t$. Such a solution exists because $I_\mu$ is continuous
and strictly increasing from zero to infinity on $[\mu,\infty)$.
Thus $\{I_\mu(M)>t\}=\{M>u_t\}$, and the Chernoff bound above gives
\[
 \Pp\{I_\mu(M)>t\}
 \le\Pp\{M\ge u_t\}
 \le e^{-I_\mu(u_t)}=e^{-t}.
\]
For the second assertion, use
$e^{z/2}-1=\frac12\int_0^z e^{t/2}\,dt$ for $z\ge0$.
Then
\begin{align*}
 \E e^{I_\mu(M)/2}
 &=1+\frac12\int_0^\infty e^{t/2}
             \Pp\{I_\mu(M)>t\}\,dt\\
 &\le1+\frac12\int_0^\infty e^{-t/2}\,dt=2.
 \qedhere
\end{align*}
\end{proof}

For a box $Q\subseteq[n]^k$, define its edge count by
\begin{equation}\label{eq:edge-count}
 e(Q)=\sum_{\mathbf i\in Q}t_{\mathbf i}.
\end{equation}
Then $\E e(Q)=p|Q|$.
Fix an integer $L\ge1$. In mode $j$, a
\emph{labeled partial partition} consists of pairwise disjoint nonempty
sets $D_j^\ell\subseteq[n]$, for a subset of labels
$\ell\in\{0,\ldots,L-1\}$. 
Each set $D_j^\ell$ is called a \emph{coordinate class} in mode $j$:
its elements are possible values of the index $i_j$.
The class is identified by the pair $(j,\ell)$; classes in different
modes are distinct even if their underlying subsets of $[n]$ coincide.
For a tuple of assigned labels $\mathbf{b}=(\ell_1,\ldots,\ell_k)$,
write $\mathbf{b}_j=\ell_j$ for its $j$th label and define
\[
 Q_{\mathbf{b}}=D_1^{\ell_1}\times\cdots\times D_k^{\ell_k},\qquad
 \mu_{\mathbf{b}}=p|Q_{\mathbf{b}}|,\qquad
 I_{\mathbf{b}}=I_{\mu_{\mathbf{b}}}(e(Q_{\mathbf{b}})).
\]
Here $\mu_{\mathbf{b}}=\E e(Q_{\mathbf{b}})$ is the expected edge count,
and $I_{\mathbf{b}}$ is the upper-tail rate \eqref{eq:rate} of the observed
count $e(Q_{\mathbf{b}})$ relative to $\mu_{\mathbf{b}}$.
Given a collection $\mathcal B$ of box labels, let
\[
 \mathcal B_j=\{\mathbf{b}_j:\ \mathbf{b}\in\mathcal B\}
\]
be the set of labels used in mode $j$. Thus $\ell\in\mathcal B_j$
exactly when $D_j^\ell$ is the $j$th factor of at least one box
in the collection. Write
\[
 h(S)=|S|\log\frac{en}{|S|},\qquad
H(\mathcal B)=\sum_{j=1}^k\sum_{\ell\in\mathcal B_j}h(D_j^\ell).
\]


Here $h(S)$ is the entropy cost of choosing a set $S$. Indeed, for $s=|S|$,
\[
 \binom ns\le\left(\frac{en}{s}\right)^s=e^{h(S)}.
\]
And
$H(\mathcal B)$ is the total entropy cost of all coordinate classes
used by the boxes in $\mathcal B$, with each class counted once even
if it occurs in several boxes. For example, when $k=2$, the boxes $D_1^0\times D_2^0$ and
$D_1^0\times D_2^1$ share the first-mode class $D_1^0$.
Thus this shared factor contributes $h(D_1^0)$, not $2h(D_1^0)$,
to $H(\mathcal B)$.


The following lemma bounds the total upper-tail rate by
$H(\mathcal B)$ and an additional term for choosing the box labels.

\begin{lemma}[Multibox discrepancy]\label{lem:local}
Fix an integer $L\ge1$ and $R>0$, and put $K=2(R+k+3)$.
With probability at least
$1-n^{-R}$, simultaneously for every labeled partial partition in every
mode and every nonempty collection $\mathcal B$ of its box labels,
\begin{equation}\label{eq:local-budget}
 \sum_{\mathbf{b}\in\mathcal B}I_{\mathbf{b}}
 \le K\bigl(H(\mathcal B)+|\mathcal B|\log(eL)\bigr).
\end{equation}
\end{lemma}

\smallskip\noindent\emph{Comparison with graph discrepancy.}
For $k=2$, the graph discrepancy estimate used in the
Kahn--Szemer\'edi argument~\cite{FO05,LR15} can be written in the
following form. Simultaneously for all nonempty $S_1,S_2\subseteq[n]$,
with $Q=S_1\times S_2$ and reference mean $\mu=p|Q|$, either
\[
 e(Q)\le C\mu
 \quad\text{or}\quad
 e(Q)\log\frac{e(Q)}{\mu}
 \le C\bigl(h(S_1)+h(S_2)\bigr).
\]
When $e(Q)\ge e^2\mu$, the quantity $e(Q)\log(e(Q)/\mu)$ is comparable
to $I_\mu(e(Q))$, so both estimates use the same large-deviation rate
in this range. The classical bound is already simultaneous over all
rectangles, but controls each rectangle separately.

In contrast, \eqref{eq:local-budget} controls the sum of the rates over
every subcollection of boxes from a coordinate partition. Adding
separate rectangle bounds repeats the entropy cost of a row or column
class whenever it is reused. The joint bound counts each shared class
only once in $H(\mathcal B)$, with an additional
$|\mathcal B|\log(eL)$ term for choosing the box labels.
This shared-class accounting is used directly in the multilevel
heavy-tuple summation below.

\begin{proof}
\emph{Fixed boxes.}
Fix a collection $\mathcal B\subseteq\{0,\ldots,L-1\}^k$ with $m\ge1$ elements.
An assignment $(D_j^\ell)_{1\le j\le k,\,\ell\in\mathcal B_j}$ is valid
if its sets are nonempty subsets of $[n]$ and pairwise disjoint within each mode.
For a fixed valid assignment, distinct label tuples give disjoint boxes,
so the counts $e(Q_{\mathbf{b}})\sim\operatorname{Bin}(|Q_{\mathbf{b}}|,p)$ are independent.
Lemma~\ref{lem:rate} therefore gives
\[
 \E\exp\left\{\frac12\sum_{\mathbf{b}\in\mathcal B}I_{\mathbf{b}}\right\}
 =\prod_{\mathbf{b}\in\mathcal B}\E e^{I_{\mathbf{b}}/2}\le2^m.
\]
 Applying Markov's
inequality to the nonnegative exponential above gives, for any $a>0$,
\begin{align*}
 \Pp\left\{\sum_{\mathbf{b}\in\mathcal B}I_{\mathbf{b}}>a\right\}
 &=\Pp\left\{\exp\left(\frac12\sum_{\mathbf{b}\in\mathcal B}I_{\mathbf{b}}\right)
                    >e^{a/2}\right\}\\
 &\le e^{-a/2}\E\exp\left(\frac12\sum_{\mathbf{b}\in\mathcal B}I_{\mathbf{b}}\right)
 \le2^m e^{-a/2}.
\end{align*}
Taking $a=K(H(\mathcal B)+m\log(eL))$ and writing
$e^{-Km\log(eL)/2}=(eL)^{-Km/2}$ yields
\begin{equation}\label{eq:fixed-budget-tail}
 \Pp\left\{\sum_{\mathbf{b}\in\mathcal B}I_{\mathbf{b}}>
       K\bigl(H(\mathcal B)+m\log(eL)\bigr)\right\}
 \le2^m(eL)^{-Km/2}e^{-KH(\mathcal B)/2}.
\end{equation}


\smallskip\noindent\emph{Uniform bound.}
To sum over coordinate sets, put
\[
 Z=\sum_{\varnothing\ne S\subseteq[n]}e^{-K h(S)/2}.
\]
Using $\binom ns\le(en/s)^s$, the monotonicity of $s\log(en/s)$ on
$[1,n]$, and $K/2-1=R+k+2$, we obtain
\begin{align}
 Z
 &=\sum_{s=1}^n\binom ns
       \exp\left\{-\frac K2s\log\frac{en}{s}\right\}\notag\\
 &\le\sum_{s=1}^n(en/s)^{-(K/2-1)s}
 \le n(en)^{-(K/2-1)}
 \le n^{-R-1}\le1.
 \label{eq:Z}
\end{align}
For fixed $\mathcal B$, the sum below chooses one nonempty subset
$D_j^\ell\subseteq[n]$ for each $j\in[k]$ and $\ell\in\mathcal B_j$,
with the sets for each fixed $j$ pairwise disjoint. Only labels in
$\mathcal B_j$ are assigned sets. In each summand, $H(\mathcal B)$
is computed from the chosen sets, and its definition gives
\[
 e^{-K H(\mathcal B)/2}
 =\prod_{j=1}^k\prod_{\ell\in\mathcal B_j}e^{-K h(D_j^\ell)/2}.
\]
All summands are nonnegative, so dropping the disjointness restrictions
can only increase the sum. Each $D_j^\ell$ can then range freely over
all nonempty subsets of $[n]$. By the distributive law for finite sums,
the unrestricted sum of these products equals the product of the
one-set sums. Hence
\begin{align*}
 \sum_{\substack{\varnothing\ne D_j^\ell\subseteq[n]\ (j\in[k],\,\ell\in\mathcal B_j)\\
                  D_j^\ell\cap D_j^{\ell'}=\varnothing\ (j\in[k],\,\ell,\ell'\in\mathcal B_j,\,\ell\ne\ell')}}
       e^{-K H(\mathcal B)/2}
 &\le\prod_{j=1}^k\prod_{\ell\in\mathcal B_j}
       \left(\sum_{\varnothing\ne S\subseteq[n]}e^{-K h(S)/2}\right)\\
 &=Z^{\sum_{j=1}^k|\mathcal B_j|}\le Z.
\end{align*}
Each factor is $Z$, and there are $\sum_j|\mathcal B_j|$ factors.
Since $\mathcal B$ is nonempty, this number is at least one;
together with $0\le Z\le1$, this proves the last inequality.

Let $\mathcal F$ be the event that there exist a nonempty collection
$\mathcal B\subseteq\{0,\ldots,L-1\}^k$ and a valid assignment
$(D_j^\ell)$ such that
\[
 \sum_{\mathbf{b}\in\mathcal B}I_{\mathbf{b}}>
 K\bigl(H(\mathcal B)+|\mathcal B|\log(eL)\bigr).
\]
By \eqref{eq:fixed-budget-tail}, the union bound, and the preceding
estimate for the sum over valid assignments,
\begin{align*}
 \Pp(\mathcal F)
 &\le\sum_{m=1}^{L^k}
     \sum_{\substack{\mathcal B\subseteq\{0,\ldots,L-1\}^k\\|\mathcal B|=m}}
       2^m(eL)^{-Km/2}Z\\
 &=Z\sum_{m=1}^{L^k}\binom{L^k}{m}2^m(eL)^{-Km/2}.
\end{align*}
The equality counts the $\binom{L^k}{m}$ label collections of size $m$.
Using $\binom{L^k}{m}\le L^{km}$ and extending the nonnegative sum gives
\[
 \Pp(\mathcal F)
 \le Z\sum_{m=1}^\infty
       \left(2e^{-K/2}L^{k-K/2}\right)^m.
\]
Since $K/2\ge k+3$ and $L\ge1$, the ratio satisfies
\[
 2e^{-K/2}L^{k-K/2}\le2e^{-(k+3)}<\frac12.
\]
The geometric series therefore yields
\[
 \Pp(\mathcal F)\le Z\sum_{m=1}^\infty2^{-m}=Z\le n^{-R}.
\]
This proves \eqref{eq:local-budget} simultaneously for all choices.
\end{proof}

\subsubsection{The deterministic heavy-tuple bound}\label{sec:heavy-sum}

We next bound the nonnegative heavy sum under the bounded degree
and multibox discrepancy assumptions. The estimate holds uniformly
over all unit vectors. 

For unit vectors $x_1,\ldots,x_k$, write $x=(x_1,\ldots,x_k)$ and
define the heavy index set at threshold $\tau=\sqrt d/n$ by
\[
 \mathcal H(x)=\left\{\mathbf i\in[n]^k:
                  \prod_{j=1}^k|x_{j,i_j}|>\tau\right\}.
\]
Thus $\mathcal H(x)=[n]^k\setminus\mathcal L(x)$.

\begin{proposition}[Deterministic multilevel summation]\label{prop:heavy}
Fix $c,\kappa,K>0$. Suppose $1\le d<n^2$, $d\ge c\log n$, and
$\tau=\sqrt d/n$. Let
\begin{equation}\label{eq:L}
 L=1+\left\lfloor\log_2(1/\tau)\right\rfloor.
\end{equation}
Let $T$ be a tensor with nonnegative entries.
Suppose that $T$ satisfies both of the following properties:
\begin{enumerate}[label=(\roman*),leftmargin=*]
\item every one-dimensional fiber has degree at most $\kappa d$;
\item for every labeled partial partition with labels in $\{0,\ldots,L-1\}$
      in each mode, and every nonempty subcollection of its boxes, the local
      entropy bound \eqref{eq:local-budget} holds with constant $K$,
      using $I_{\mathbf{b}}=I_{\mu_{\mathbf{b}}}(e(Q_{\mathbf{b}}))$ and $\mu_{\mathbf{b}}=(d/n)|Q_{\mathbf{b}}|$.
\end{enumerate}
Then there is a constant $C$, depending only on $k,c,\kappa,K$, such that
\[
 \sup_{\|x_1\|_2=\cdots=\|x_k\|_2=1}
 \sum_{\mathbf i\in\mathcal H(x)}
       t_{\mathbf i}\prod_{j=1}^k|x_{j,i_j}|
 \le C\sqrt d.
\]
\end{proposition}
\begin{proof}
Fix a tensor $T$ satisfying \textup{(i)}--\textup{(ii)} and unit vectors
$x_1,\ldots,x_k$. We bound the heavy sum uniformly in these vectors.

\subsubsection*{Dyadic decomposition}
Put $a_\ell=2^{-\ell}$ and, for $0\le\ell<L$, let
\[
 D_j^\ell=\{i:\ |x_{j,i}|>\tau,\quad a_\ell/2<|x_{j,i}|\le a_\ell\}.
\]
Since $2^{-L}<\tau\le2^{-(L-1)}$, these classes partition the indices
with $|x_{j,i}|>\tau$. Retain their labels when discarding empty
classes to obtain a labeled partial partition. Define
\begin{equation}\label{eq:alpha}
 \alpha_{j,\ell}=|D_j^\ell|a_\ell^2,\qquad
 \sum_{\ell=0}^{L-1}\alpha_{j,\ell}
 \le4\sum_{\ell=0}^{L-1}\sum_{i\in D_j^\ell}|x_{j,i}|^2
 \le4.
\end{equation}
Unless stated otherwise, level sums below range over $0,\ldots,L-1$;
empty classes have mass zero.

For each nonempty box $Q_{\mathbf b}=D_1^{\ell_1}\times\cdots\times D_k^{\ell_k}$,
write $a_j=a_{\ell_j}$, $s_j=|D_j^{\ell_j}|$, and
$\alpha_j=\alpha_{j,\ell_j}$. Set
\begin{equation}\label{eq:box-parameters}
 a_{\mathbf b}=\prod_j a_j,\qquad
 \mu_{\mathbf b}=\frac dn\prod_j s_j>0,\qquad
 \rho_{\mathbf b}=\frac{a_{\mathbf b}}{\tau},\qquad
 \lambda_{\mathbf b}=\frac{e(Q_{\mathbf b})}{\mu_{\mathbf b}}.
\end{equation}
Let $\mathcal B$ index the nonempty boxes with $a_{\mathbf b}>\tau$.
Every heavy tuple has $|x_{j,i_j}|>\tau$ in each mode, because all
coordinate magnitudes are at most one. It therefore belongs to a
unique dyadic box, whose amplitude satisfies
$a_{\mathbf b}\ge\prod_j|x_{j,i_j}|>\tau$. Since $T$ is nonnegative,
\begin{equation}\label{eq:heavy-box-domination}
 \sum_{\mathbf i\in\mathcal H(x)}t_{\mathbf i}\prod_j|x_{j,i_j}|
 \le\sum_{\mathbf b\in\mathcal B}a_{\mathbf b}e(Q_{\mathbf b}).
\end{equation}

Following the class-by-class organization of \cite[Section~4.3]{ZZ21},
partition the box labels into
\begin{align*}
 \mathcal C_1
 &=\{\mathbf b\in\mathcal B:\ na_j^2\le\rho_{\mathbf b}
                                      \text{ for some }j\},\\
 \mathcal C_2
 &=\{\mathbf b\in\mathcal B\setminus\mathcal C_1:
                               \lambda_{\mathbf b}\le e^2\rho_{\mathbf b}\},\\
 \mathcal C_3&=\mathcal B\setminus(\mathcal C_1\cup\mathcal C_2).
\end{align*}
We control these classes by fiber degrees, dyadic masses, and joint
discrepancy, respectively.

\subsubsection*{Boxes in $\mathcal C_1$: the fiber-degree bound}
Fix a mode $j$ and the other level labels, and put
$q=(\prod_{i\ne j}a_i)/\sqrt d$. Since
$\rho_{\mathbf b}=na_{\mathbf b}/\sqrt d$, the condition
$na_j^2\le\rho_{\mathbf b}$ is equivalent to $a_j\le q$.
As $\ell_j$ varies, the sets $D_j^{\ell_j}$ are disjoint.
Thus the corresponding boxes use disjoint portions of the same
$\prod_{i\ne j}s_i$ mode-$j$ fibers. By \textup{(i)},
\[
 \sum_{\substack{\ell_j:\ \mathbf b\in\mathcal B\\a_j\le q}}e(Q_{\mathbf b})
 \le\kappa d\prod_{i\ne j}s_i.
\]
Using $a_j\le q$ before summing gives
\[
 \sum_{\substack{\ell_j:\ \mathbf b\in\mathcal B\\a_j\le q}}
       \frac{a_{\mathbf b}e(Q_{\mathbf b})}{\sqrt d}
 \le\frac{q\prod_{i\ne j}a_i}{\sqrt d}\,
       \kappa d\prod_{i\ne j}s_i
 =\kappa\prod_{i\ne j}\alpha_i.
\]
Sum over the other labels and then over $j$. Every box in
$\mathcal C_1$ is counted at least once, and all terms are nonnegative.
Factoring the sums over independent labels and using \eqref{eq:alpha},
\begin{equation}\label{eq:unbalanced-sum}
 \sum_{\mathbf b\in\mathcal C_1}
       \frac{a_{\mathbf b}e(Q_{\mathbf b})}{\sqrt d}
 \le\kappa\sum_{j=1}^k
       \prod_{i\ne j}\left(\sum_{\ell_i}\alpha_{i,\ell_i}\right)
 \le k\kappa\,4^{k-1}.
\end{equation}

\subsubsection*{Boxes in $\mathcal C_2$: the mass bound}
Since $|Q_{\mathbf b}|a_{\mathbf b}^2=\prod_j\alpha_j$,
\[
 \frac{a_{\mathbf b}e(Q_{\mathbf b})}{\sqrt d}
 =\frac{\lambda_{\mathbf b}}{\rho_{\mathbf b}}\prod_j\alpha_j
 \le e^2\prod_j\alpha_j
 \qquad(\mathbf b\in\mathcal C_2).
\]
Enlarging the nonnegative mass sum to all label tuples and factoring
it, \eqref{eq:alpha} gives
\begin{equation}\label{eq:moderate-sum}
 \sum_{\mathbf b\in\mathcal C_2}
       \frac{a_{\mathbf b}e(Q_{\mathbf b})}{\sqrt d}
 \le e^2\prod_{j=1}^k\left(\sum_{\ell_j}\alpha_{j,\ell_j}\right)
 \le e^2 4^k.
\end{equation}

\subsubsection*{Boxes in $\mathcal C_3$: the joint discrepancy bound}
There is nothing to prove if $\mathcal C_3$ is empty. Otherwise,
its defining conditions are
\begin{equation}\label{eq:high}
 1<\rho_{\mathbf b}<na_j^2\quad(j\in[k]),\qquad
 \lambda_{\mathbf b}>e^2\rho_{\mathbf b}.
\end{equation}
Hence \eqref{eq:rate} yields
\[
 I_{\mathbf b}\ge e(Q_{\mathbf b})(\log\lambda_{\mathbf b}-1)
               \ge e(Q_{\mathbf b})\log(e\rho_{\mathbf b}).
\]
For $s\ge0$, put
\[
 \mathcal C_3(s)=\left\{\mathbf b\in\mathcal C_3:
        s<\frac{a_{\mathbf b}}{\sqrt d\log(e\rho_{\mathbf b})}\right\},
 \qquad H(\varnothing)=0.
\]
Each box belongs to this subcollection on an interval of length
$a_{\mathbf b}/(\sqrt d\log(e\rho_{\mathbf b}))$.
Tonelli's theorem and \textup{(ii)} therefore give
\begin{align}
 \sum_{\mathbf b\in\mathcal C_3}\frac{a_{\mathbf b}e(Q_{\mathbf b})}{\sqrt d}
 &\le\int_0^\infty\sum_{\mathbf b\in\mathcal C_3(s)}I_{\mathbf b}\,ds
 \notag\\
 &\le K\int_0^\infty H(\mathcal C_3(s))\,ds
       +K\log(eL)\int_0^\infty|\mathcal C_3(s)|\,ds.
 \label{eq:high-integrated}
\end{align}
Hypothesis \textup{(ii)} covers every subcollection, so it applies
even though $\mathcal C_3(s)$ depends on the box counts.

For the entropy integral, if $D_j^\ell$ occurs in $\mathcal C_3(s)$,
some box using it satisfies \eqref{eq:high}. The monotonicity of
$u/\log(eu)$ on $[1,\infty)$ gives
\[
 s<\frac{\rho_{\mathbf b}}{n\log(e\rho_{\mathbf b})}
 \le\frac{a_\ell^2}{\log(en a_\ell^2)},\qquad na_\ell^2>1.
\]
Each such class is counted once in $H(\mathcal C_3(s))$, for an
interval of length at most the displayed bound. Its integrated
entropy cost is therefore at most
\[
 \frac{h(D_j^\ell)a_\ell^2}{\log(en a_\ell^2)}
 =\frac{\alpha_{j,\ell}\log(en a_\ell^2/\alpha_{j,\ell})}
          {\log(en a_\ell^2)}
 \le\alpha_{j,\ell}\left(1+
       \frac{\log_+(1/\alpha_{j,\ell})}{\log(en a_\ell^2)}\right).
\]
Classes unused by $\mathcal C_3$ contribute zero.
Summing these bounds and applying Lemma~\ref{lem:mass} in each mode,
with $t_\ell=na_\ell^2=n4^{-\ell}$ and the masses in \eqref{eq:alpha},
yields
\begin{equation}\label{eq:total-entropy-charge}
 \int_0^\infty H(\mathcal C_3(s))\,ds
 \le\sum_{j=1}^k\sum_{\substack{0\le\ell<L\\na_\ell^2>1}}
       \alpha_{j,\ell}\left(1+
       \frac{\log_+(1/\alpha_{j,\ell})}{\log(en a_\ell^2)}\right)
 \le\frac{28}{3}k.
\end{equation}

For the second integral, $\log(e\rho_{\mathbf b})>1$ gives
\begin{equation}\label{eq:overhead-sum}
 \int_0^\infty|\mathcal C_3(s)|\,ds
 =\sum_{\mathbf b\in\mathcal C_3}
       \frac{a_{\mathbf b}}{\sqrt d\log(e\rho_{\mathbf b})}
 \le\frac1{\sqrt d}\sum_{\ell_1,\ldots,\ell_k\ge0}
       2^{-\ell_1-\cdots-\ell_k}
 =\frac{2^k}{\sqrt d}.
\end{equation}
Since $d\ge1$ and $d\ge c\log n$, \eqref{eq:L} implies
$L\le1+\log_2 n\le C_c d$, and hence
$\log(eL)/\sqrt d\le C_c$.
Substituting these estimates into \eqref{eq:high-integrated} gives
\begin{equation}\label{eq:high-sum}
 \sum_{\mathbf b\in\mathcal C_3}\frac{a_{\mathbf b}e(Q_{\mathbf b})}{\sqrt d}
 \le K\left(\frac{28}{3}k+
                 \frac{2^k\log(eL)}{\sqrt d}\right)
 \le C_{k,c,K}.
\end{equation}

Combining \eqref{eq:unbalanced-sum}, \eqref{eq:moderate-sum}, and
\eqref{eq:high-sum} in \eqref{eq:heavy-box-domination} proves the desired
bound. The constant is independent of the chosen vectors, so taking
the supremum completes the proof.
\end{proof}

\subsection{Completion of the proof}\label{sec:completion}

We combine the light estimate, the two events used for the heavy
estimate, and the net approximation \eqref{eq:net-bound}.


\begin{proof}[Proof of Theorem~\ref{thm:main}]
By increasing $C_{k,r,c}$ if necessary,
we may assume that $n$ is sufficiently large that
$d\ge c\log n\ge1$. Set $R=r+2$.
Apply Lemma~\ref{lem:local} with the deterministic $L$ in
\eqref{eq:L}, Lemma~\ref{lem:fibers}, and Lemma~\ref{lem:light}.
Their three events hold simultaneously except on an event of
probability at most
$3n^{-(r+2)}\le n^{-r}$.
On this intersection, fix unit vectors $x_1,\ldots,x_k$ and put
$q_{\mathbf i}=\prod_j|x_{j,i_j}|$.
Proposition~\ref{prop:heavy}, with $d=np$, gives
\[
 \sum_{\mathbf i\in\mathcal H(x)}t_{\mathbf i}q_{\mathbf i}
 \le C_{\mathrm h}\sqrt d,
\]
where $C_{\mathrm h}$ depends only on $k,r,c$, and the bound holds
uniformly over all such unit vectors.

The mean tensor is $P=\E T=p\,\one^{\otimes k}$, whose entries all
equal $p$, so $W=T-P$ has entries $w_{\mathbf i}=t_{\mathbf i}-p$.
The heavy contribution of $P$ in absolute value is bounded by
$p\sum_{\mathbf i\in\mathcal H(x)}q_{\mathbf i}$.
For a heavy tuple, $q_{\mathbf i}>\tau$, hence
$q_{\mathbf i}\le q_{\mathbf i}^2/\tau$. Therefore
\begin{equation}\label{eq:heavy-mean}
\begin{aligned}
 p\sum_{\mathbf i\in\mathcal H(x)}q_{\mathbf i}
 &\le\frac p\tau\sum_{\mathbf i\in\mathcal H(x)}q_{\mathbf i}^2
 \le\frac p\tau\sum_{\mathbf i\in[n]^k}q_{\mathbf i}^2\\
 &=\frac p\tau\prod_{j=1}^k\|x_j\|_2^2
 =\frac{np}{\sqrt d}=\sqrt d.
\end{aligned}
\end{equation}
Here we used $\|x_j\|_2=1$, $\tau=\sqrt d/n$, and $d=np$.

Now take $x=(x_1,\ldots,x_k)\in\mathcal T^k$.
Splitting the full multilinear form into light and heavy tuples gives
\[
 W(x_1,\ldots,x_k)
 =L_W(x)+\sum_{\mathbf i\in\mathcal H(x)}
                 (t_{\mathbf i}-p)\prod_{j=1}^k x_{j,i_j}.
\]
By the triangle inequality and $|t_{\mathbf i}-p|\le t_{\mathbf i}+p$,
\begin{align*}
 |W(x_1,\ldots,x_k)|
 &\le |L_W(x)|+\sum_{\mathbf i\in\mathcal H(x)}|t_{\mathbf i}-p|q_{\mathbf i}\\
 &\le |L_W(x)|
       +\sum_{\mathbf i\in\mathcal H(x)}t_{\mathbf i}q_{\mathbf i}
       +p\sum_{\mathbf i\in\mathcal H(x)}q_{\mathbf i}.
\end{align*}
The first term is at most $B_{k,R}\sqrt d$ by Lemma~\ref{lem:light}.
The second is at most $C_{\mathrm h}\sqrt d$ by the heavy bound above,
and the third is at most $\sqrt d$ by \eqref{eq:heavy-mean}. Consequently,
\[
 \max_{x\in\mathcal T^k}|W(x_1,\ldots,x_k)|
 \le (B_{k,R}+C_{\mathrm h}+1)\sqrt d.
\]

We now apply \eqref{eq:net-bound} to the full multilinear form $W$,
after the light and heavy contributions have been combined. This gives
\[
 \|W\|\le
       2\max_{(y_1,\ldots,y_k)\in\mathcal T^k}
            |W(y_1,\ldots,y_k)|\le C_{k,r,c}\sqrt d.
\]
This proves \eqref{eq:main}.
\end{proof}


\begin{proof}[Proof of Corollary~\ref{cor:weighted}]
If $A=0$, the result is immediate. By homogeneity, assume
$\|A\|_\infty=1$.
Pad $A$ with zeros to an $N\times\cdots\times N$ tensor and extend
$\Omega$ by zero entries, with the corresponding probabilities also
set to zero. Since the added entries
of $A$ are zero, this does not change either random tensor in the
statement or the spectral and entrywise maximum norms of $A$. It is
therefore enough to consider equal dimensions. 
Set
$d=Np_*$, $\tau=\sqrt d/N$, and $R=r+2$.
By increasing $C_{k,r,c}$ if necessary, we may assume that
$d\ge c\log N\ge1$.

We use the proof of Theorem~\ref{thm:main} with the centered entries
$a_{\mathbf i}(\omega_{\mathbf i}-p_{\mathbf i})$. For fixed vectors
in the product net, the light summands are independent, centered, and
bounded in absolute value by $\tau$. Their total variance is at most
\[
 \sum_{\mathbf i}p_{\mathbf i}a_{\mathbf i}^2
        \prod_{j=1}^k x_{j,i_j}^2
 \le p_*\prod_{j=1}^k\|x_j\|_2^2=p_*.
\]
Thus the Bernstein estimate and the union bound in
Lemma~\ref{lem:light} hold with  probability at least $1-N^{-R}$.

For the two inputs to Proposition~\ref{prop:heavy}, use independent
uniform random variables $U_{\mathbf i}$ on $[0,1]$ to realize
\[
 \omega_{\mathbf i}=\one\{U_{\mathbf i}\le p_{\mathbf i}\},
 \qquad
 \omega^*_{\mathbf i}=\one\{U_{\mathbf i}\le p_*\}.
\]
Then $\Omega^*=(\omega^*_{\mathbf i})$ has independent
Bernoulli$(p_*)$ entries, and $\Omega\le\Omega^*$ entrywise.
Thus every fiber degree and box count of $\Omega$ is bounded by its
counterpart for $\Omega^*$. Since $u\mapsto I_{p_*|Q|}(u)$ is
nondecreasing, the simultaneous degree and discrepancy bounds for
$\Omega^*$ in Lemmas~\ref{lem:fibers} and~\ref{lem:local} also hold
for $\Omega$, with reference means $p_*|Q|=(d/N)|Q|$ and total failure
probability at most $2N^{-R}$.

The light, fiber-degree, and discrepancy events therefore hold
simultaneously with probability at least
$1-3N^{-R}\ge1-N^{-r}$. Since $p_*\le1$, we have
$d\le N<N^2$, so Proposition~\ref{prop:heavy} applies to $\Omega$
on this event.

It remains to compare the heavy part with the nonnegative count tensor
$\Omega$. Write $q_{\mathbf i}=\prod_j|x_{j,i_j}|$. Since
$|a_{\mathbf i}|\le1$,
\[
 \left|\sum_{\mathbf i\in\mathcal H(x)}
       a_{\mathbf i}(\omega_{\mathbf i}-p_{\mathbf i})
       \prod_jx_{j,i_j}\right|
 \le\sum_{\mathbf i\in\mathcal H(x)}
       (\omega_{\mathbf i}+p_{\mathbf i})q_{\mathbf i}.
\]
Proposition~\ref{prop:heavy} bounds the part containing
$\omega_{\mathbf i}$. For the mean part,
\[
 \sum_{\mathbf i\in\mathcal H(x)}p_{\mathbf i}q_{\mathbf i}
 \le\frac{p_*}{\tau}\sum_{\mathbf i}q_{\mathbf i}^2
 =\frac{p_*}{\tau}=\sqrt d.
\]
Both bounds are uniform over the unit vectors. Combining the light and
heavy estimates and applying the same product-net argument proves
\eqref{eq:weighted}.
\end{proof}

\section{Proof of Theorem~\ref{thm:fw}}\label{sec:fw-proof}

Throughout this section, $X,d,P$ are as in \eqref{eq:fw-model}.
We use the deterministic heavy bound from Proposition~\ref{prop:heavy}.
Its probabilistic inputs are a fiber-degree estimate and a discrepancy
bound for dependent box counts. The light contribution is a sum over
the independent edge draws.

\subsection{A joint rate estimate for multinomial counts}

The counts in disjoint boxes, together with the count outside their
union, have a multinomial distribution. A sampled edge belongs to at
most one box. Consequently, for nonnegative parameters, their joint
moment generating function is bounded by that of independent Poisson
variables with the reference means below. We use this to prove the
joint rate estimate needed in Lemma~\ref{lem:local}. The reference
means may exceed the actual means; in our application they are equal.

\begin{lemma}[Joint exponential moment]\label{lem:multinomial}
Let $Q_1,\ldots,Q_s$ be pairwise disjoint nonempty subsets of $[n]^k$,
and write $e(Q_b)=\sum_{\mathbf i\in Q_b}X_{\mathbf i}$ for their
occupancy counts under \eqref{eq:fw-model}.
Choose arbitrary reference means
\[
 \mu_b\ge \frac{m|Q_b|}{n^k}>0,
 \qquad I_b=I_{\mu_b}(e(Q_b)).
\]
Then
\begin{equation}\label{eq:mixed-mgf}
 \E\exp\left(\frac12\sum_{b=1}^s I_b\right)\le2^s.
\end{equation}
\end{lemma}
\begin{proof}
Put $q_b=|Q_b|/n^k$. Disjointness, $1+x\le e^x$, and
$mq_b\le\mu_b$ give, for $\theta_1,\ldots,\theta_s\ge0$,
\begin{align*}
 \E\exp\left(\sum_{b=1}^s\theta_be(Q_b)\right)
 &=\left(1+\sum_{b=1}^s q_b(e^{\theta_b}-1)\right)^m\le\exp\left(\sum_{b=1}^s\mu_b(e^{\theta_b}-1)\right).
\end{align*}
For $B\subseteq[s]$ and $t_b\ge0$, choose $u_b\ge\mu_b$ with
$I_{\mu_b}(u_b)=t_b$, as in Lemma~\ref{lem:rate}.
Since $\{I_b>t_b\}=\{e(Q_b)>u_b\}$, Markov's inequality with
$\theta_b=\log(u_b/\mu_b)$ on $B$ and zero elsewhere yields
\[
 \Pp\{I_b>t_b\text{ for every }b\in B\}
 \le \exp\left(\sum_{b\in B}
       [\mu_b(e^{\theta_b}-1)-\theta_bu_b]\right)
 =e^{-\sum_{b\in B}t_b}.
\]
Finally, expand $\prod_{b=1}^s e^{I_b/2}$ using
$e^{z/2}-1=\frac12\int_0^\infty e^{t/2}\one\{z>t\}\,dt$ for $z\ge0$.
Tonelli's theorem and the joint tail bound give
\begin{align*}
 \E\exp\left(\frac12\sum_{b=1}^s I_b\right)
 &\le\sum_{B\subseteq[s]}\prod_{b\in B}
 \left(\frac12\int_0^\infty e^{-t/2}\,dt\right)
 =\sum_{B\subseteq[s]}1=2^s.
 \qedhere
\end{align*}
\end{proof}

\begin{corollary}[Multinomial multibox discrepancy]\label{cor:multi-local}
Fix a deterministic integer $L\ge1$ and $R>0$.
For the model \eqref{eq:fw-model}, define, for each nonempty box
$Q_{\mathbf{b}}$ of a fixed labeled partial partition,
\[
 e(Q_{\mathbf{b}})=\sum_{\mathbf i\in Q_{\mathbf{b}}}X_{\mathbf i},
 \qquad \mu_{\mathbf{b}}=\E e(Q_{\mathbf{b}})=\frac dn|Q_{\mathbf{b}}|,
 \qquad I_{\mathbf{b}}=I_{\mu_{\mathbf{b}}}(e(Q_{\mathbf{b}})).
\]
Here $e(Q_{\mathbf{b}})$ counts sampled edges in $Q_{\mathbf{b}}$ with
multiplicity, and $I_{\mathbf{b}}$ measures its upper-tail deviation
from the mean $\mu_{\mathbf{b}}$ using the rate function
\eqref{eq:rate}. In particular, $I_{\mathbf{b}}=0$ when
$e(Q_{\mathbf{b}})\le\mu_{\mathbf{b}}$.
With probability at least $1-n^{-R}$, every partial
partition and every nonempty subcollection of its boxes satisfy
\[
 \sum_{\mathbf{b}\in\mathcal B}I_{\mathbf{b}}
 \le K\bigl(H(\mathcal B)+|\mathcal B|\log(eL)\bigr),
 \qquad K=2(R+k+3).
\]
\end{corollary}
\begin{proof}
For each fixed assignment of coordinate classes, the boxes are
disjoint and their means are $(d/n)|Q_{\mathbf{b}}|$.
Lemma~\ref{lem:multinomial} replaces the independence step in
\eqref{eq:fixed-budget-tail}, giving exactly the same bound.
The counting of used labels, coordinate sets, and box collections
in \eqref{eq:Z} and the remainder of Lemma~\ref{lem:local} is
deterministic and applies without change. Summing those bounds
gives $n^{-R}$.
\end{proof}

\subsection{Fiber degrees and light tuples for sampled edges}

For each fixed fiber its degree is
$\operatorname{Bin}(m,n^{-(k-1)})$, with mean $d$.
The Chernoff and union-bound calculation in
Lemma~\ref{lem:fibers} therefore proves the same degree estimate.

For the light part, put $W=X-P$, fix unit vectors and any threshold
$\tau>0$, and define
\[
 c_{\mathbf i}=\prod_{j=1}^k x_{j,i_j}\,
                   \one\left\{\left|\prod_{j=1}^k x_{j,i_j}\right|
                                    \le\tau\right\}.
\]
The centered light sum is
\begin{equation}\label{eq:sample-light}
 L_W(x):=\sum_{\mathbf i}c_{\mathbf i}(X_{\mathbf i}-d/n)
 =\sum_{\nu=1}^m\left(c_{E_\nu}-\E c_{E_\nu}\right).
\end{equation}
The summands on the right are independent and bounded in absolute value by
$2\tau$, and their total variance is at most
\begin{equation}\label{eq:sample-variance}
 m\,\E c_{E_1}^2
 =\frac{m}{n^k}\sum_{\mathbf i}c_{\mathbf i}^2
 \le\frac dn.
\end{equation}
With $\tau=\sqrt d/n$, Bernstein's inequality~\cite[Theorem~2.9.5]{V26}
gives
\begin{equation}\label{eq:sample-light-tail}
 \Pp\{|L_W(x)|>B\sqrt d\}
 \le2\exp\left\{-n\frac{B^2}{2+4B/3}\right\}.
\end{equation}
For $R>0$, choose $B=B_{k,R}$ with
$B^2/(2+4B/3)\ge k\log(1+4k)+R+1$.
Since the product net $\mathcal T^k$ from Section~\ref{sec:discretization}
has at most $(1+4k)^{kn}$ points, a union bound in
\eqref{eq:sample-light-tail} gives, for $n\ge2$,
\[
 \Pp\left\{\max_{x\in\mathcal T^k}|L_W(x)|>B\sqrt d\right\}
 \le2\exp\left\{n\left[k\log(1+4k)-\frac{B^2}{2+4B/3}\right]\right\}
 \le n^{-R}.
\]

\subsection{Completion of the proof}

\begin{proof}[Proof of Theorem~\ref{thm:fw}]
By increasing $C_{k,r,c}$ if necessary, we may assume that $n$
is sufficiently large that $d\ge c\log n\ge1$.

Suppose first that $d<n^2$.
Take $R=r+2$ in
Corollary~\ref{cor:multi-local}, with $L$ as in \eqref{eq:L}.
The corollary, the fiber bound, and
\eqref{eq:sample-light-tail} followed by a product-net union bound
hold together with probability at least $1-3n^{-R}$.
Proposition~\ref{prop:heavy} controls the heavy part on this event.
The contribution of the mean $P$ on heavy tuples is at most
$\sqrt d$, by the calculation in \eqref{eq:heavy-mean} with $p=d/n$.
The net approximation \eqref{eq:net-bound} now gives the conclusion.

Finally, if $d\ge n^2$, then $\tau=\sqrt d/n\ge1$.
Every coordinate product of unit vectors is at most one in absolute
value, so every tuple is light. In this case
\eqref{eq:sample-light-tail} and the same net approximation suffice.
Increasing the constant over these cases proves the theorem.
\end{proof}

\appendix
\section{A dyadic entropy estimate}\label{app:dyadic}

The estimate for boxes in $\mathcal C_3$ uses the following bound
for the entropy sum in one mode. It is uniform in the number of
dyadic levels.

\begin{lemma}[Entropy of dyadic masses]\label{lem:mass}
Let $n\ge2$ and $L\ge1$ be integers, and set
\[
 t_\ell=n4^{-\ell},\qquad \ell=0,\ldots,L-1.
\]
If $\alpha_\ell\ge0$ and $\sum_{\ell=0}^{L-1}\alpha_\ell\le4$, then
\begin{equation}\label{eq:mass-entropy}
 \sum_{\substack{0\le\ell<L\\t_\ell>1}}\alpha_\ell
 \left(1+\frac{\log_+(1/\alpha_\ell)}{\log(et_\ell)}\right)
 \le 8+\frac43.
\end{equation}
Terms with $\alpha_\ell=0$ are interpreted as zero.
\end{lemma}

\begin{thebibliography}{99}
\bibitem[]{ACS25} Al-Ghattas, O.; Chen, J.; Sanz-Alonso, D.. Sharp concentration of simple random tensors. Information and Inference: A Journal of the IMA 14 (2025) iaaf029
\bibitem[]{BFSU98} Broder, A. Z.; Frieze, A. M.; Suen, S.; Upfal, E.. Optimal construction of edge-disjoint paths in random graphs. SIAM Journal on Computing 28 (1998) 541--573
\bibitem[]{BGJLR24} Bandeira, A. S.; Gopi, S.; Jiang, H.; Lucca, K.; Rothvoss, T.. A geometric perspective on the injective norm of sums of random tensors. (2024)
\bibitem[]{BKW22} Bamberger, S.; Krahmer, F.; Ward, R.. The Hanson--Wright inequality for random tensors. Sampling Theory, Signal Processing, and Data Analysis 20 (2022) 14
\bibitem[]{Boe24} Boedihardjo, M. T.. Injective norm of random tensors with independent entries. (2024)
\bibitem[]{CGJ18} Cook, N.; Goldstein, L.; Johnson, T.. Size biased couplings and the spectral gap for random regular graphs. The Annals of Probability 46 (2018) 72--125
\bibitem[]{CS26} Chen, J.; Sanz-Alonso, D.. Sharp concentration of simple random tensors II: asymmetry. Information and Inference: A Journal of the IMA 15 (2026) iaag010
\bibitem[]{DJPP13} Dumitriu, I.; Johnson, T.; Pal, S.; Paquette, E.. Functional limit theorems for random regular graphs. Probability Theory and Related Fields 156 (2013) 921--975
\bibitem[]{DM24} Dartois, S.; McKenna, B.. Injective norm of real and complex random tensors I: From spin glasses to geometric entanglement. (2024)
\bibitem[]{DM26} Dartois, S.; McKenna, B.. A moment-based approach to the injective norm of random tensors. (2026)
\bibitem[]{DWZ25} Dumitriu, I.; Wang, H.-X.; Zhu, Y.. Partial recovery and weak consistency in the non-uniform hypergraph stochastic block model. Combinatorics, Probability and Computing 34 (2025) 1--51
\bibitem[]{FKS89} Friedman, J.; Kahn, J.; Szemer{\'e}di, E.. On the second eigenvalue of random regular graphs. Proceedings of the Twenty-First Annual ACM Symposium on Theory of Computing (STOC '89) (1989) 587--598
\bibitem[]{FO05} Feige, U.; Ofek, E.. Spectral techniques applied to sparse random graphs. Random Structures \& Algorithms 27 (2005) 251--275
\bibitem[]{FW95} Friedman, J.; Wigderson, A.. On the second eigenvalue of hypergraphs. Combinatorica 15 (1995) 43--65
\bibitem[]{GJL25} Guan, J.; Jiang, B.; Li, Z.. On decomposability and subdifferential of the tensor nuclear norm. (2025)
\bibitem[]{JO14} Jain, P.; Oh, S.. Provable tensor factorization with missing data. Advances in Neural Information Processing Systems 27 (2014) 1431--1439
\bibitem[]{LCL20} Lei, J.; Chen, K.; Lynch, B.. Consistent community detection in multi-layer network data. Biometrika 107 (2020) 61--73
\bibitem[]{LLV17} Le, C. M.; Levina, E.; Vershynin, R.. Concentration and regularization of random graphs. Random Structures \& Algorithms 51 (2017) 538--561
\bibitem[]{LMSZ26} Luo, H.; Ma, A.; Stephan, L.; Zhu, Y.. Wedge sampling: Efficient tensor completion with nearly-linear sample complexity. (2026)
\bibitem[]{LP26} Lucca, K.; Pesenti, L.. Norm bounds for sparse random tensors and spectral gap of random hypergraphs. (2026)
\bibitem[]{LR15} Lei, J.; Rinaldo, A.. Consistency of spectral clustering in stochastic block models. The Annals of Statistics 43 (2015) 215--237
\bibitem[]{NDT15} Nguyen, N. H.; Drineas, P.; Tran, T. D.. Tensor sparsification via a bound on the spectral norm of random tensors. Information and Inference: A Journal of the IMA 4 (2015) 195--229
\bibitem[]{TS14} Tomioka, R.; Suzuki, T.. Spectral norm of random tensors. (2014)
\bibitem[]{TY19} Tikhomirov, K.; Youssef, P.. The spectral gap of dense random regular graphs. The Annals of Probability 47 (2019) 362--419
\bibitem[]{V20} Vershynin, R.. Concentration inequalities for random tensors. Bernoulli 26 (2020) 3139--3162
\bibitem[]{V26} Vershynin, R.. High-Dimensional Probability: An Introduction with Applications in Data Science. (2026)
\bibitem[]{XY21} Xia, D.; Yuan, M.. Effective tensor sketching via sparsification. IEEE Transactions on Information Theory 67 (2021) 1356--1369
\bibitem[]{YZ17} Yuan, M.; Zhang, C.-H.. Incoherent tensor norms and their applications in higher order tensor completion. IEEE Transactions on Information Theory 63 (2017) 6753--6766
\bibitem[]{Z23} Zhu, Y.. On the second eigenvalue of random bipartite biregular graphs. Journal of Theoretical Probability 36 (2023) 1269--1303
\bibitem[]{ZZ21} Zhou, Z.; Zhu, Y.. Sparse random tensors: Concentration, regularization and applications. Electronic Journal of Statistics 15 (2021) 2483--2516
\end{thebibliography}
\end{document}
